\documentclass[a4paper]{extarticle}
\usepackage[fontsize=14pt]{scrextend}
\usepackage[
margin=0.25in
]{geometry}
\usepackage{setspace}
%\setstretch{1.4}
\usepackage{xcolor, multicol}
\pagestyle{empty}
\usepackage[inline]{enumitem}
\setlist{itemjoin=\hfil, label=\hspace*{10pt}\arabic*)\hspace{10pt}, labelindent=0pt, labelsep=0pt}
\usepackage{tikz}
\setlength{\parindent}{0pt}
\setlength{\columnseprule}{0.4pt}
\setlength{\columnsep}{1cm}
\usepackage{amsmath, amssymb, mathtools, amsthm, nccmath}

\usepackage[main=arabic, bidi=basic, provide=*, layout=counters.tabular graphics]{babel}

%\babelprovide[justification=kashida, transforms=kashida.base]{arabic}
%\babelfont{rm}[Script=Arabic, Scale=1.33]{Scheherazade}
\babelfont{rm}[Script=Arabic]{Times New Roman}
\babelfont[english]{rm}{Times New Roman}

\directlua{Babel.arabic.justify_factor = 1.03}

\usepackage[lite]{mtpro2}
\DeclareMathSymbol{0}{\mathalpha}{operators}{`0}
\DeclareMathSymbol{1}{\mathalpha}{operators}{`1}
\DeclareMathSymbol{2}{\mathalpha}{operators}{`2}
\DeclareMathSymbol{3}{\mathalpha}{operators}{`3}
\DeclareMathSymbol{4}{\mathalpha}{operators}{`4}
\DeclareMathSymbol{5}{\mathalpha}{operators}{`5}
\DeclareMathSymbol{6}{\mathalpha}{operators}{`6}
\DeclareMathSymbol{7}{\mathalpha}{operators}{`7}
\DeclareMathSymbol{8}{\mathalpha}{operators}{`8}
\DeclareMathSymbol{9}{\mathalpha}{operators}{`9}

\def\SPACE{\hspace*{28pt}}
\def\VSPACE{\vskip5pt}
\def\MARKS{\hfill\textbf{(10 درجات)}}
\def\EMPTY{\text{---------}}

\newcommand{\dashedrule}{%
	\noindent\leaders\hbox{\rule{5pt}{0.4pt}\hskip3pt}\hfill\kern0pt
	\vskip10pt
}

\AtBeginDocument{
	\allowdisplaybreaks
	\renewcommand{\jot}{10pt}
	\abovedisplayskip=7pt
	\belowdisplayskip=7pt
}

\begin{document}


%\noindent
\begin{minipage}{0.3\textwidth}
	\begin{center}
		\textbf{ثانوية آمنة بنت وهب}\linebreak
		\textbf{للبنات}
	\end{center}
\end{minipage}
\hfill
\begin{minipage}{0.3\textwidth}
	\begin{center}
		\textbf{اسئلة امتحان نصف السنة}\\
		\textbf{للعام الدراسي}
		\textbf{2025 - 2026}\\
		\textbf{(الدور الاول)}
	\end{center}
\end{minipage}
\hfill
\begin{minipage}{0.3\textwidth}
	\begin{center}
		\textbf{المادة : الرياضيات}\\
		\textbf{الصف : الاول المتوسط}\\
		\textbf{التأريخ : 24/1/2026}\\
		\textbf{الوقت : ساعة ونصف}
	\end{center}
\end{minipage}
{\vskip5pt\hrule\vskip5pt
\centering
	\textbf{ملاحظة : الاجابة عن خمسة اسئلة فقط (لكل سؤال 20 درجة)}
\vskip5pt\hrule\vskip5pt}

\textbf{س 1: أ/}
	 جدي ناتج ما يأتي لأثنين فقط \hfill \textbf{(10 درجات)}
\\\SPACE 
\begin{enumerate*}
	\item $\dfrac{7}{9} - \dfrac{8}{5}$
	\item $4.6+(-3.7)$
	\item $1\dfrac{2}{3} \times \dfrac{-3}{10}$
\end{enumerate*}\\

\textbf{س 1: ب/}
	 جدي ناتج الضرب $4\times(8y+4z+5)$ \hfill \textbf{(10 درجات)}

\dashedrule

\textbf{س 2:}
	 اجيبي عن فرعين فقط\\
\SPACE \textbf{أ/} 
	قدّر الجذر (لأثنين فقط) 
\begin{enumerate*}
\item $\sqrt[3]{24}$
\item $\sqrt{13}$
\item $\sqrt{94}$
\end{enumerate*}\hfill\textbf{(10 درجات)}
\VSPACE
\SPACE
\textbf{ب/}
	 من $24z^2wy$ اطرح $-32z^2wy$ \hfill\textbf{(10 درجات)}
\VSPACE
\SPACE\textbf{ج/} طلى عصام غرفة بلون جديد خلال 12 ساعة فأذا ساعده رياض بطلاء غرفة اخرى بالقياس نفسه ، الى كم ساعة يجتاجون 
\SPACE\hspace{10pt}
لأنجاز العمل ؟ (بأعتبار ان رياض ينجز العمل بكفاءة عصام وسرعته)
\hfill\textbf{(10 درجات)}

\dashedrule

\textbf{س 3: أ/} مثلث قائم الزاوية طولا ضلعيه (12 ، 5 سم) جد طول الوتر \hfill\textbf{(10 درجات)}
\VSPACE
\SPACE\textbf{ب/} جدي القيمة العددية للمقادير الجبرية التالية (لواحد فقط) \MARKS\\\SPACE\SPACE
	\begin{enumerate*}[itemjoin=]
		\item $3r^2 + 2v + 16, \quad r = 3, v = 7$\\
		\SPACE\SPACE\SPACE
		\item $8w - 7z + 12, \quad w = 3, z = 3$
	\end{enumerate*}
	
\dashedrule

\textbf{س 4: أ/} شريط طوله $3\dfrac{1}{5} \,m$ قطع الى 4 قطع متساوية ما طول القطعة الواحدة ؟\MARKS
\VSPACE
\SPACE
\textbf{ب/} اجيبي عن نقطة واحدة \MARKS\\
\SPACE\SPACE \textbf{1)} حل معادلة الجمع والطرح باستعمال العلاقة بين الجمع والطرح 
$X + 22 = | -42 | $

\SPACE\SPACE\textbf{2)} حل معادلة القسمة باستعمال العلاقة بين الضرب والقسمة 
$42 \div Y = -6$

\dashedrule

\textbf{س 5: اجيبي عن فرعين فقط}\\
\SPACE\textbf{أ/} استعمل ترتيب العمليات وجد ناتج
$(32- 9) \times (18 \div 6) + 4$ \MARKS
\VSPACE\SPACE
\textbf{ب/} اذا كانت
$
A = \{-4, -3, -1, 0, 2\}
$
،\SPACE
$
B = \{-4, -1, 3, 5\}
$\MARKS\\\SPACE\SPACE
جدي 
\begin{enumerate*}[itemjoin=\hspace{15pt}]
	\item $A\cap B$ \item $A\cup B$
\end{enumerate*}

\VSPACE\SPACE
\textbf{ج/} جدي التقسيم التناسبي $3:4$ من 210000 \MARKS

\dashedrule

\textbf{س 6: أ/} اكملي العبارات التالية بما يناسبها \MARKS
\begin{enumerate}[topsep=0pt, itemsep=0pt, leftmargin=60pt]
	\item $8^0 = \EMPTY$
	\item $\sqrt[3]{125} = \EMPTY$
	\item النسبة المئوية للعدد $\mfrac{20}{5}$ هي \EMPTY
	\item الحد الجبري $7xy$ المعامل هو \EMPTY\, والمتغير هو \EMPTY
\end{enumerate}

\VSPACE\VSPACE\SPACE
\textbf{ب/} استعملي خصائص العمليات لتحسبي
 $(21 + 33) + 9$ \MARKS
 
 
 \hrule
 
 \begin{center}
 	\textbf{مدرسة المادة : رنا سمير عبدالعباس}
 \end{center}
 
\end{document}



